\begin{papercard}{hruschka2004}
\cf{Stages} \stagetag{C} \stagetag{I}. An \emph{empirical} account of the iterative codebook-revision loop in three CDC HIV studies. It shows that two coders agree poorly on the first draft codebook but reach high $\kappa$ after a finite number of rounds, identifies factors that speed this up, and gives eight guidelines (appendix) that extend MacQueen et al.

\cf{Position and claim} Intercoder reliability is contested (positivist critique vs.\ skeptics vs.\ applied team-based work). The authors take the third view: reliability is a necessary, not sufficient, quality check on one stage of the research chain \pc{hruschka2004}{4}{ensuring reliability is a necessary (but not sufficient) step for drawing accurate conclusions}. Three findings: coders initially produce very different codings; agreement improves through codebook revision; simple procedural steps speed this up \pc{hruschka2004}{3}{we show that (1) coders initially generate very different codings of text}.

\cf{Process (Figure 1)} Segment the text, create the codebook, code, assess reliability, modify the codebook, final coding. Codebook creation: team members independently propose themes, then agree on a master list, with inclusion and exclusion rules per code. The simplest codes are dichotomous \pc{hruschka2004}{6}{Although it is possible to have codes with multiple values (e.g., high, medium, low), the simplest codes are dichotomous}. A lead coder distributes a random subset (e.g.\ 60 of 300, or all of 30) and each member codes independently, then the team meets. Statistic: Cohen's $\kappa$ rather than percent agreement \pc{hruschka2004}{6}{can dramatically overestimate the true degree of intercoder reliability by not taking chance agreement into account}. Final dataset: both coders code the full set, intermediate checks (e.g.\ at 50\,\%), and per-code final $\kappa$ \pc{hruschka2004}{9}{reliability checks may be made at intermediate stages of this final coding}.

\cf{Interpretation frameworks} Landis \& Koch (0.81--1.00 almost perfect), Cicchetti (0.75--1.00 excellent; stricter, clinical) and Miles \& Huberman ($\approx$0.90). The studies used stringent cut-offs of $\kappa\ge0.80$ or $0.90$ \pc{hruschka2004}{7}{In the studies presented below, we used fairly stringent cutoffs at kappa}.

\cf{The three studies (Table 1)} Hemophilia (295 coded interviews, 70--80 codes, 4 questions, researcher notes), Zimbabwe (70 respondents, 200+ codes across 27 questions, 4 analysed), Los Angeles bathhouse (24 interviews, 330--340 codes, 30 questions, transcripts of one line to a page). Two coders each; coding was dichotomous per code and segment; reports came from CDC EZ-Text. Cut-offs: hemophilia $\kappa>0.8$; LA 80\,\% of codes with $\kappa>0.9$; Zimbabwe 90\,\% of codes with $\kappa>0.9$ \pc{hruschka2004}{10}{The LA bathhouse study required that 80\% of codes have a kappa score greater than 0.9}.

\cf{Results} Round 1 was always poor: 32.9\,\% of codes reached $\kappa\ge.8$ (hemophilia), 39.0\,\% reached $\ge.9$ (LA), and 38.5--61.1\,\% (Zimbabwe questions) \pc{hruschka2004}{12}{Regardless of study or question, the first round of qualitative coding generated low levels of intercoder reliability}. Rounds needed: 2 (hemophilia, final 64.6\,\%), 8 (LA), 3--9 (Zimbabwe, depending on the question) \pc{hruschka2004}{13}{one question (Q6A) required nine rounds, while two questions required four rounds and one required three rounds}. Faster convergence with fewer codes per round (about 30+ codes: nine rounds; $<$20 codes: three to four) \pc{hruschka2004}{13}{Furthermore, the number of rounds required to achieve acceptable levels of intercoder reliability was reduced from nine}. Long unsegmented text units (up to a page) hurt even when codes were few. Other factors: interviewer quality and building the first codebook \emph{after} reading the texts \pc{hruschka2004}{14}{coders from the LA bathhouse study noted that the initial codebook was generated prior to viewing any of the interviews}.

\cf{Interpretive convergence} Improvement may come from clearer definitions \emph{or} from coders' repeated discussions converging (Hak \& Bernts), so the codebook may only work for that team. Test: a second team applies the final codebook independently \pc{hruschka2004}{15}{The possibility of interpretive convergence arises most obviously in the LA bathhouse study}. Other limits: only two coders per study, no validity assessment of the codebook, short semistructured data.

\cf{Appendix guidelines} (1) Reliability will not fix poor design or data. (2) Organise texts, codebooks and a single codebook editor. (3) \emph{Code independent subsamples each round} so agreement is not a memory of earlier discussion. (4) Limit the codes per round to under 20 and code one question at a time; question-specific codes beat global codes \pc{hruschka2004}{19}{we suggest a further reduction to fewer than twenty codes, and coding only one question at a time}. (5) Do not sacrifice relevance for reliability (the trivial code ``THE'') \pc{hruschka2004}{19}{Because of the simplicity of this recognition task, two coders would likely be able to generate very similar assignments}. (6) Never use percent agreement alone. (7) Choose the cut-off by purpose: $\ge.7$--$.8$ for rough prevalence estimates, $\ge.9$ for individual-level classification \pc{hruschka2004}{20}{If one is only concerned with a rough estimate of a code’s population prevalence, a cutoff of kappa}. (8) Report clearly, with code-specific $\kappa$, and interpret low-agreement codes cautiously \pc{hruschka2004}{21}{the code-specific kappa should be included because it may be very different from the summary measure}.

\cf{Use for this thesis} Evidence for the pilot design. (1) Expect poor first-round agreement; plan 2--4 rounds, with a fresh random subset each time and a stop rule stated beforehand (e.g.\ share of codes above the threshold). (2) Keep codes per pass under about 20: split the codebook into families (privacy policy, consent, controls, risk communication) and code one family at a time. (3) Build the first codebook after reading a sample of real walkthrough output. (4) Use a threshold justified by purpose; the descriptive prevalence of barriers allows a lower cut-off than classification of individual apps. (5) Report per-code $\kappa$ and name \emph{interpretive convergence} as a limitation. $\kappa$ itself has prevalence problems, which Hruschka notes only in passing (see Feinstein \& Cicchetti, Gwet).
\end{papercard}
